% ── 다리 초안: dreyer2010/dreyer2011 → 우리 spinodal gap + (iii-a) 평탄역 ──
% 주 삽입 위치: ch1_sec07_broadening.tex, (iii-a) 문단(현행 68–74행) 직후,
%   "(iii-b) broadening 층..."(현행 76행) 앞.
% 앵커 문장: "...거시적 \emph{평탄역}이다(삽입 전지 다입자 평형의 표준 그림; Dreyer 등\cite{dreyer2010,dreyer2011})."
% 부속(선택) 삽입: ch1_sec04_hys.tex 현행 17행 직후 한 줄 포인터(중복 회피 위해 본 다리로 통합 권장).
% 컨테이너: srcbox. 우리 기호로 서술.
% ★확인 필요: 원 논문의 정확한 닫힌 수식·식번호는 미확인(웹 PDF 파싱 불가). 방법 요지는
%   Nature Mater. 9,448 (2010)·Contin. Mech. Thermodyn. 23,211 (2011) 초록·본문 구조로 검증 —
%   대응은 "결과 구조" 수준(우리 식 ↔ 논문 결과)으로 제시, 논문 식 재현은 하지 않음.

\begin{srcbox}
\textbf{다리 --- 이 평탄역과 §4 gap 이 어느 다입자 결과인가.} (iii-a) 의 공통 전위 순차 전환과 §4 의 분기
gap(식~\eqref{eq:dUhys})은 같은 Dreyer 다입자 열역학\cite{dreyer2010,dreyer2011} 의 두 얼굴이다. 대응은
결과 구조 수준이다(왼쪽 본문 기호, 오른쪽 원 논문 결과):
\begin{center}\footnotesize
\begin{tabular}{@{}ll@{}}
\hline
본문 & Dreyer\cite{dreyer2010,dreyer2011} 결과(개념) \\
\hline
단일입자 비단조 $V_{\eq,j}(\xi)$ (식~\eqref{eq:Veq}, $\Omega_j{>}2RT$) & 비단조 단일입자 화학퍼텐셜/등온선 \\
과주행 상한 $\xi_{s,j}^\pm$ (식~\eqref{eq:spinodal}) & 준안정 가지의 안정성 한계 \\
방전$\cdot$충전 분기 $\Delta U_j^\hys$ (식~\eqref{eq:dUhys}) & 영전류에도 남는 충$\cdot$방전 전위 갈림 \\
(iii-a) 공통 전위 순차 전환 평탄역 & 입자별 순차 전환의 겉보기 평형(apparent equilibrium) \\
\hline
\end{tabular}
\end{center}
\emph{방법 요지} --- Dreyer 등은 다수 입자가 빠른 리튬 교환으로 \emph{하나의 공통 화학퍼텐셜}을 공유할 때,
단일입자 비단조 등온선 위에서 입자들이 동시에 두-상으로 가지 않고 \emph{순차}로 전환해 겉보기 평형 평탄역을
만들며, 방전$\cdot$충전이 서로 다른 준안정 가지를 타 영전류에서도 갈림이 남음(히스테리시스)을 다입자
열역학으로 보인다. \emph{가정 차} --- 원 골격은 입자 간 교환이 충분히 빨라 전 입자가 공통 $\mu$ 를 공유하고
(평균장$\cdot$size-blind) 단일입자 등온선이 갈림의 유일 원천이나, 본문은 그 위에 겉보기 전위
$U_\app{=}U_j{+}\eta$ 의 앙상블 이질성(\S\ref{sec:broadening} (iii-b) 의 $\eta$ 분포, 비-크기
heterogeneity)을 별도 broadening 층으로 얹고 평형 중심 $U_j$ 는 입자무관 상수로 두며, gap 은 우리 spinodal
상한 식~\eqref{eq:dUhys} 으로 닫아 원 논문의 실측 갈림과 \emph{부호}만 대조한다(식~\eqref{eq:dUhys} 유도의
$\Omega_j{>}2RT$ 문턱은 본문 고유).
\end{srcbox}
